% DRAFT, not wired into paper/main.tex or thesis/. Written 2026-11-09, extended 2026-11-10 (sections 4 to 6) for Sans's review. Nothing here is final.
% Every number is tagged in a comment with its source. Items marked TODO need a source or a result that does not exist yet.
% Scope reminders: Gemma-4-E2B-it only, BFCL v4 multi_turn_base, one Mac; training history stays in the appendix; statistics wording is never stronger than the record.

%% ---------------------------------------------------------------------------------------------
%% (1) Chapter 6, one paragraph to add after the self-distillation paragraph
%% ---------------------------------------------------------------------------------------------
\paragraph{A larger model as teacher did not change the outcome.} We also trained the model on tool-call answers written by a much larger model (DeepSeek V4.1 Flash, called through a commercial API), keeping only answers that the task checker accepted, in several rounds with different data mixes and training objectives. None of the resulting adapters (small trained add-on weights; Appendix~\ref{app:train}) beat the untuned model on the tasks we held out for screening. The last round scored 59 and 64 of 155 screening tasks for the two trained variants, against 68 and 67 for two runs of the untuned model, and a follow-up set of diagnostic runs could not separate any trained or random-direction adapter from the untuned model (Appendix~\ref{app:train}). % SOURCE: docs/DISTILLATION_ROUND5_FINDINGS.md (A2 table); Round 6 stage 1 below. TODO(Sans): confirm the chapter keeps "several rounds" or names the count.

%% ---------------------------------------------------------------------------------------------
%% (2) Appendix, training history: new paragraphs
%% ---------------------------------------------------------------------------------------------
\paragraph{Teacher-student rounds.} The teacher was DeepSeek V4.1 Flash, the student Gemma-4-E2B-it, served as in the rest of this thesis. Tasks were written by the teacher in the BFCL style (a starting state, several user turns, and the calls that should answer each turn), checked by running them in the benchmark's own backend, and solved by the teacher inside the same tool loop the student uses; only checker-verified answers became training rows. The screening set (call it the screening bank) held 155 tasks, 81 of them from an earlier round and 74 from the last; neither part is training data, and the first part is not independent of our choices because its aggregate results helped decide what to try next. % SOURCE: docs/DISTILLATION_ROUND5_FINDINGS.md "Screening bank and rule".
The pre-registered screening rule asked a trained variant to reach at least 77 of 155 and to beat the better untuned run (68) before it earned a repeat. It was a way to decide where to spend compute, not a significance test. % SOURCE: same section.

\paragraph{Last round, supervised training.} Two adapters were trained on identical rows (315 training rows from 95 tasks, 80 optimizer updates, rank-8 adapter on the query projections of the last 16 layers, 524{,}288 trainable parameters): one with the standard loss and one that weighted the answer tokens more than the reasoning text. They scored 59 and 64 of 155 screening tasks (29 and 33 of 81; 30 and 31 of 74), against 68 and 67 for two untuned runs. Neither reached the rule, and both were below both untuned runs. The 5-task gap between the two trained variants is descriptive only. Two untuned runs on the 81-task part exchanged 5 gains and 5 losses with equal totals, so differences of a few tasks between single runs are noise. % SOURCE: Round 5 findings, A2 table and integrity notes.
A second arm (preference training on pairs of the student's own good and bad answers at the same state) was never trained: the pre-registered floor was 150 training pairs from at least 50 tasks, and the pipeline produced 120 pairs from 60 tasks. % SOURCE: Round 5 findings, section B.

\paragraph{Diagnostic runs (Round 6, stage 1).} To ask whether the training update itself or the training content hurt, we served the last-round supervised adapter at full, half and quarter strength, a random-direction adapter with the same size and singular-value spectrum, and the untuned model, on the same 155 screening tasks (greedy decoding, same server build and settings, 0 errors in every run). The composites were 64 (adapter, full strength), 66 (half), 65 (quarter), 64 (random direction), and 64, 65 and 68 for three untuned runs taken in the same session. % SOURCE: THESIS_PROGRESS.md 2026-11-09 evening, stage 1 and extra base run; attempts 20261109_121007_stage1.1ihURW and 20261109_191001_extrabase.zf0OhZ in the Mac worktree Thesis-wt-r6. TODO(Sans): the run-name table in the appendix needs these attempt ids.
No arm can be separated from the untuned model in this session, and the untuned model itself moved by 4 tasks between runs. We therefore make no statement about whether the earlier adapter was harmful or harmless. Two of the registered consistency checks flagged the session (the first untuned run, 64, fell below its expected range of 65 to 70, and the adapter at full strength, 64, rose above its expected 56 to 62), and one additional untuned run was added as registered. % SOURCE: Codex review docs/reviews/codex_s1_stage1_result.md; r6_tools summarize flags.

\paragraph{Data quality, not only data amount.} A separate measurement asked how many usable tasks the teacher pipeline yields. From 2{,}000 authoring slots, 1{,}819 were attempted before a time limit; 181 produced a valid task, 99 remained after the teacher-solve and held-out filters, and the cost was \$2.11 over 13{,}587 calls. A blind accuracy review of 40 of the 99 tasks (does the reference answer do what the user asked, not only does it run) found 7 wrong tasks by one reviewer model and 3 by two other reviewers; an added state check plus a request-coverage judge removed 28 of the 99 tasks (71 remained), and a fresh 40-task review of tasks that passed the added checks found 1 wrong by one reviewer and 0 by two others. % SOURCE: Part 1 YIELD_REPORT.json; docs/reviews/codex_a1_sample_verdicts.md; DATA_YIELD_FACTORY_FIX.md calibration; docs/reviews/codex_a1b_fresh_sample_verdicts.md and subagent_a1b_fresh_sample_verdicts.md in the Mac worktree Thesis-wt-data.
The defects found before the added checks were of a kind that execution-based checking cannot see: a requested action missing from the reference, a quoted value (``the card ending in 7712'') matching nothing in the state, and new objects whose id silently overwrote an existing one. % SOURCE: Codex A1 verdicts, verified by the assistant against the sample text.

%% ---------------------------------------------------------------------------------------------
%% (3) Not yet writable; waiting on registered results
%% ---------------------------------------------------------------------------------------------
% G1 result paragraph (source: Thesis-wt-r6 sft/data/distill/r6g1/attempts/20261109_201851_g1.P8Hx7x/G1_REPORT.json; docs/TRAINING_RESULTS_LOG.md section 3; Codex review "No major problems found").
\paragraph{An overfit check: the training pipeline can learn.} Before scaling the data, we asked whether our training and serving pipeline can learn anything at all. We trained a wider adapter (layers 8 to 34, rank 16) for four passes on 243 rows from 40 synthetic tasks, with the model's reasoning mode switched off, and then scored the same 40 tasks. The untrained model solved 11 of the 40 tasks; the adapter solved 32. The rule we fixed beforehand (at least 24 of 40 and at least 8 more than the untrained model) was met. At the level of single steps the adapter reproduced 234 of the 243 training answers exactly, against 165 for the untrained model. This is a pipeline check, not a result about generalization: the tasks used for scoring were the training tasks, the run used the earlier tool-list format rather than the compact list, four tasks lost rows to a length limit, and it is one run with one seed. It says nothing about the adapters of the earlier rounds.
% TODO(G2/G3, G2 authoring done and judge pending, G3 base references running): scaled data and a trained comparison against the compact-list stack (Run 84, 88/200 = 44.0%; same stack with the original list, Run 85, 72/200 = 36.0%). A BFCL-200 claim needs Sans's explicit go and the paired test already used in Chapter 5; the existing statistics hedges (suggestive, not significant after Holm correction) stay.

%% ---------------------------------------------------------------------------------------------
%% (3b) Rounds 2 to 4, filled 2026-11-10 from the findings documents in the Mac worktree Thesis-wt-distill (read-only). Internal run names stay out of the text.
%% ---------------------------------------------------------------------------------------------
\paragraph{Earlier teacher-student rounds.} Before the last round the same kind of training was tried three times on small screening banks (Table~\ref{tab:distill-rounds}). The first round trained on 154 teacher-written rows and on the model's own replies as a control, and compared the results on a 20-task check and on a 23-task bank of fresh teacher-verified tasks. The second trained two adapters (one on 320 rows built from 95 tasks mixing teacher-informed and own-reply rows, one on 304 rows of the model's own verified replies from 43 tasks) and compared them on a 103-task bank. The third studied error recovery: a harness nudge that tells the model to retry after a failed call, and recovery training rows. In every case the untuned model scored at least as many tasks as each trained variant, and the registered rules stopped each round before a larger test.
% SOURCES (all in Thesis-wt-distill/docs): DISTILLATION_ROUND2_FINDINGS.md (mid20 table; dev bank table; "Headline"), DISTILLATION_ROUND3_FINDINGS.md (result table; reading; addendum), DISTILLATION_ROUND4_FINDINGS.md (Stage A table; Stage B). The paragraph's row counts: 154 teacher rows (Round 2, "t1"), 320 rows / 95 tasks (AB) and 304 rows / 43 tasks (A) from the Round 3 addendum and reading sections.

\begin{table}[!htbp]
\centering
\caption{Teacher-student screening rounds. Each score is passed tasks out of the bank size shown; ``untuned'' is the same model served the same way with no adapter. Greedy decoding, one run per variant unless two untuned runs are shown. A pass here is a screening result, not an accuracy estimate.}
\label{tab:distill-rounds}
\begin{tabular}{@{}p{2.6cm}p{3.0cm}p{2.1cm}p{3.2cm}p{3.3cm}@{}}
\toprule
Round & What was varied & Bank & Untuned & Trained or changed \\
\midrule
First & Rows written by the teacher (154) vs the model's own replies & 20-task check; 23 fresh tasks & 20-task check: 17, 15, 16 in three runs; fresh tasks: 7 & 20-task check: 10 (teacher rows), 13--14 (own-reply control); fresh tasks: 6, 4, 5, 4 for four variants \\
Second & Two adapters, 320 rows vs 304 rows & 103 tasks & 44 and 44 & 41 and 36 (registered bar: 47) \\
Third & Retry nudge in the harness; recovery rows & 81 tasks & 35 and 35 & Nudge: 35; earlier adapter: 30. Recovery rows: only 45 strict rows from 166 candidates, below the registered floor of 100; not trained \\
Fourth & Weighted answer loss vs standard loss & 155 tasks & 68 and 67 & 64 and 59 (registered bar: 77 and above 68) \\
\bottomrule
\end{tabular}
\end{table}
% TABLE SOURCES. First round: ROUND2 findings, "mid20 controls" (Runs 69 = 17, 90 = 15, 88 = 16; t1 = 10; c1 = 13 fused, 14 runtime) and "dev bank" (7/23 untuned; t1 6, t2 4, t3 5, c1 4). Second round: ROUND3 findings result table (B1 44, B2 44, AB 41, A 36; bar max(B1,B2)+3 = 47). Third round: ROUND4 findings Stage A table (B1 35, B2 35, H 35, AB 30) and Stage B (166 candidates, 60 pass the implemented filter, 45 strict; rule < 100 stops). Fourth round: ROUND5 findings A2 table. NAMING: the findings documents number the rounds 2, 3, 4, 5; the table above uses First/Second/Third/Fourth for the rounds that precede the Round 6 diagnostics, TODO(Sans): decide whether the thesis names them by number.
% HEDGES TO KEEP (from the same documents): the nominal untuned configuration scored 17 and then 15 of 20 on a rerun (environment equivalence not established); two untuned runs on the 81-task bank exchanged 5 gains and 5 losses; the Round 3 training prompts lacked the thinking marker that serving renders (all 659 rows differed from the served prompt; adding the marker made all of them identical), so Round 3 is confounded and should not be read as a clean test of that recipe; teacher-forced validation loss fell in every round (for example 0.270 to 0.193) without a gain in passes; the Round 3 and Round 4 training rows were filtered by whether the student could already reach the action.
% COSTS (logged API cost, per-call usage logs, not an account balance): Round 2 about $0.70 (about $1.55 over all calls to that point), Round 3 about $2.18 of a $4 cap, Round 4 about $3.2 of a $15 allowance, Round 5 $1.4453 of a $2.5 cap.
% LIMITS to state in the final text: single seed per variant; one teacher; screening tasks drawn from the same authoring pipeline as the training tasks; the earlier bank is adaptively informed; Round 6 diagnostics are single runs with session drift flagged.

%% ---------------------------------------------------------------------------------------------
%% (4) Data pipeline for the scaled run. Added 2026-11-10 from docs/TRAINING_RESULTS_LOG.md (sections 2, 4 to 8e); every count there names its file. Internal run and attempt names stay out of the text.
%% ---------------------------------------------------------------------------------------------
\paragraph{Scaling up the training data.} After the diagnostic runs, the question became whether the earlier adapters failed because they were trained on too few and too uneven tasks. We therefore rebuilt the data pipeline to produce several hundred tasks that had been checked for correctness, and trained on them (next paragraphs). A task here is a synthetic multi-turn tool-calling conversation written by the teacher model (a starting state of the benchmark's simulated services, several user requests, and the calls that should answer each one). Four checks stand between an authored task and a training row. First, the task must run without error in the benchmark's own backend. Second, deterministic state checks reject tasks whose reference answer cannot be what the user asked for (for example a newly created object whose identifier silently overwrites an existing one). Third, a language-model judge reads the user requests and the reference calls and rejects tasks where a requested action is missing or an invented one is present. Fourth, the teacher must solve the task itself, in the same tool loop and the same compact tool-list format that the student sees at test time, and only solutions the benchmark checker accepts become rows. Tasks sharing long word sequences with the test tasks or with the screening banks are removed at the end, and the result is checked again row by row (provenance check). % SOURCE: results log sections 2, 4, 6; docs/DATA_YIELD_FACTORY_FIX.md; docs/DATA_YIELD_G2_PREREGISTRATION.md.
\paragraph{Yield of the pipeline.} Table~\ref{tab:distill-funnel} gives the counts for the two authoring batches. The second batch used more varied themes, a prompt rule against the identifier overwrite, and a slot mix chosen to cover service combinations that the first batch lacked. It yielded fewer tasks per slot, mostly because the two service combinations it favored (messaging and travel booking) produced valid tasks in only about 2 to 3\,\% of slots in both batches, and it ended early: one authoring process reached its cost cap and the spend monitor stopped the other two in their last chunk, so 938 slots (668 in the two stopped chunks, 270 refused by the cost guard) produced nothing and the authoring step is reported as stopped by its cap rather than as passed. The 270 refusals made by the cost guard after that cap were counted as lost slots, not as provider errors; there were no provider errors in this batch. In the first batch, 189 of 193 provider errors came from one process and were later traced to the account key reaching its monthly limit, not to the provider failing. Together the two batches give 501 tasks. % SOURCE: results log 4, 4b, 8, 8b, 8c; G2B prereg sections 5, 6. NUMBERS: 938 = 668 slots in the two killed chunks + 270 refused. TODO(Sans): decide whether the thesis names the cap stop in the main text or leaves it to the appendix table of attempts.
\begin{table}[!htbp]
\centering
\caption{Task counts through the data pipeline. A slot is one authoring attempt by the teacher model. ``Valid'' means the task ran in the benchmark backend and its reference answer was accepted. Batch~1 counts after the held-out filter are reduced by one task removed by the provenance check.}
\label{tab:distill-funnel}
\begin{tabular}{@{}lrr@{}}
\toprule
Step & Batch 1 & Batch 2 \\
\midrule
Slots planned & 10{,}302 & 10{,}002 \\
Valid tasks authored & 824 & 393 \\
After de-duplication against earlier tasks & 734 & 368 \\
Pass state checks and judge & 481 & 270 \\
At least one teacher solution accepted & 344 & 195 \\
After held-out filter (final tasks) & 325 & 177 \\
After provenance check & 324 & 177 \\
\midrule
Training rows, direct-answer format (train + validation) & 2{,}667 + 272 & 1{,}416 + 125 \\
Training rows, reasoning-first format (train + validation) & 2{,}494 + 249 & 1{,}309 + 117 \\
\bottomrule
\end{tabular}
\end{table}
% TABLE SOURCES: results log 4, 4b, 6, 8, 8b, 8c: attempt reports (author/report.json, merge_report.json, prefilter_report.json, solve/pilot_report.json, G2_REPORT.json). Batch 1 ``planned'' is the number of slots worked (10,302); batch 2 planned 10,002, worked 9,334. Batch 2 reasoning-first rows are after removing 4 rows (3 train, 1 validation) in the provenance check. Ledger cost: batch 1 \$10.56, batch 2 \$10.51, Part 1 \$2.12; the API key's own counter reads \$31.68 and the difference is not explained (results log 8c) -- state the ledger figures only, with that caveat, or leave costs to the appendix.
\paragraph{Checking the tasks themselves.} Because execution-based checks cannot tell whether a reference answer does what the user asked, we had the tasks read by other models. In the first pilot, 40 sampled tasks were read by three reviewers (one coding model and two general models); the coding model marked 7 wrong and the other two 3, all of which it had also marked. That sample failed the rule we had set beforehand (at most 2 wrong per reviewer), so we added the state checks and the judge described above, which removed 28 of the 99 pilot tasks. A fresh sample of 40 tasks that passed the new checks had 1 wrong (coding model) and 0 (the others) and met the rule. We also tried a mechanical repair for the identifier overwrite instead of discarding those tasks: it kept 129 of 134 affected tasks, but a 20-task review found 6 wrong, including cases where the repair moved the identifier that the reference answer relied on. That failed the rule and the repaired tasks were not used. For the second batch, a 40-task sample of the final tasks had 2 wrong (coding model) and 0 and 0 (the others), which meets the rule at its limit; both wrong tasks concern a simulator convention (the ``sent messages'' call in the benchmark's messaging service returns the whole inbox), not a mistake in the task, and nine of the 501 tasks contain that call. In 46 of the 501 tasks a trading order is cancelled by identifier without first reading its details, which reviewers called ambiguous; we leave these tasks in. The reviewers are models, not people, and each sample has 40 tasks. % SOURCE: results log sections 2, 4 (repair), 8d; Codex = the coding-model reviewer. NOTE FOR SANS: the advisor's request listed three passed samples; the record has two passed (fresh pilot sample, batch 2) and two failed (first pilot sample, repair sample). The batch 1 final tasks have no post-solve sample review on record (results log section 6 says so). Do not state three until it exists.

%% ---------------------------------------------------------------------------------------------
%% (5) G3 design (arms A and B), threshold fixed before scoring. Results are placeholders.
%% ---------------------------------------------------------------------------------------------
\paragraph{Design of the scaled training run.} The scaled run trains the wider adapter used in the overfit check (layers 8 to 34, rank 16) for two passes over the batch-1 rows, in two formats. In the direct-answer format (arm A) the target is the list of calls with the model's reasoning mode switched off. In the reasoning-first format (arm B) the target is two or three plain sentences of reasoning followed by the same calls, with the reasoning mode switched on; the sentences were written by the teacher model and were scanned for leaked phrases and for overlap with the test material. Arm B has 6.5\,\% fewer training rows than arm A because rows with a missing, one-sentence or four-sentence reasoning text, or a leaked phrase, were excluded. The adapter is served on the untuned model's exact weights as a run-time add-on, and both arms are scored on the same 155-task screening bank, in the same compact tool-list format, with the reasoning mode matching the arm. % SOURCE: docs/DISTILLATION_G3_PREREGISTRATION.md sections 1 to 5; results log 5 to 7. TODO(Sans): the 6.5 percent is 1 - 2494/2667 = 6.5 percent [ANALYSIS in the results log reads 93.5 percent of A].
\paragraph{Scoring rule fixed in advance.} Before any trained model was scored we ran the untuned model on the screening bank three times in the new format: with reasoning on it solved 59 and 61 of 155 tasks and with reasoning off 54. A trained arm earns a repeat run only if its first run solves at least 70 of 155 (the mean of the two reasoning-on runs, 60.0, plus 6\,\% of 155 rounded up, and more than the better of the two); a repeat must reach the same bar. The rule is a screening device, not a significance test. No result from the 200 test tasks enters this step. % SOURCE: results log 5 (G3_BASE_REPORT.json: on1 59, on2 61, off 54; threshold 70).
\paragraph{Results of the scaled run.} % PLACEHOLDER. Arm A (direct answers): run 1 = TODO of 155, repeat = TODO. Arm B (reasoning first): run 1 = TODO of 155, repeat = TODO. Source files: attempts G3_A_RUN1.json, G3_B_RUN1.json in Thesis-wt-r6 sft/data/distill/r6g3/attempts. The text must state the pass or fail against 70, the caveat that the screening bank shares its author with the training data, and that only the 200 test tasks settle a claim.
% Exploratory probe (results log 8e): arm A after the first of two passes solved 37 of 81 tasks on the first screening part (reasoning off); the untuned references on that part are 30 and 33 (reasoning on) and 27 (off). Descriptive only; a mid-schedule checkpoint; do not put in the main text without the caveat. 
\paragraph{Combined retrain.} % PLACEHOLDER. One further run on all 501 tasks in the arm that scored higher on the screening bank (ties to arm A). Because the arm is selected on that bank, the result is described as exploratory. Pre-registration: docs/DISTILLATION_G3C_COMBINED_PREREGISTRATION.md. Result: TODO.
% Hedges that stay with any result paragraph: single seed; screening banks are not independent of the authoring pipeline; session drift of the untuned model (4 tasks between runs in the diagnostic session); validation loss is logged and never used to choose anything; the 200-task claim needs the paired test of Chapter 5 and Sans's explicit go.
